\documentclass[12pt,a4paper]{report}

%%% Hlavní soubor. Zde se definují základní parametry a~odkazuje se na ostatní části. %%%

%% Verze pro jednostranný tisk:
% Okraje: levý 40mm,~pravý 25mm,~horní a~dolní 25mm
% (ale pozor,~LaTeX si sám přidává 1in)
\setlength\textwidth{145mm}
\setlength\textheight{247mm}
\setlength\oddsidemargin{15mm}
\setlength\evensidemargin{15mm}
\setlength\topmargin{0mm}
\setlength\headsep{0mm}
\setlength\headheight{0mm}
% \openright zařídí,~aby následující text začínal na pravé straně knihy
\let\openright=\clearpage

%% Pokud tiskneme oboustranně:
% \documentclass[12pt,a4paper,twoside,openright]{report}
% \setlength\textwidth{145mm}
% \setlength\textheight{247mm}
% \setlength\oddsidemargin{14.2mm}
% \setlength\evensidemargin{0mm}
% \setlength\topmargin{0mm}
% \setlength\headsep{0mm}
% \setlength\headheight{0mm}
% \let\openright=\cleardoublepage

\usepackage[pdfa,colorlinks=false,urlcolor=black,bookmarksopen=true]{hyperref} 	% Záložky v~souboru

%% Vytváříme PDF/A-2u
\usepackage[a-2u]{pdfx}

%% Přepneme na českou sazbu a~fonty Latin Modern
\usepackage[czech]{babel}
\usepackage{lmodern}
\usepackage[T1]{fontenc}
\usepackage{textcomp}
\usepackage{imakeidx}		% Index
%% Nastavení indexu
\makeindex[columns=2,title=Index]

%% Použité kódování znaků: obvykle latin2,~cp1250 nebo utf8:
\usepackage[utf8]{inputenc}
\usepackage{csquotes}
\MakeOuterQuote{"}

%%% Další užitečné balíčky (jsou součástí běžných distribucí LaTeXu)
\usepackage{amsmath}        % rozšíření pro sazbu matematiky
\usepackage{amsfonts}       % matematické fonty
\usepackage{amsthm}         % sazba vět,~definic apod.
\usepackage{amssymb}
\usepackage{bbding}         % balíček s~nejrůznějšími symboly
			    % (čtverečky,~hvězdičky,~tužtičky,~nůžtičky,~...)
\usepackage{bm}             % tučné symboly (příkaz \bm)
\usepackage{graphicx}       % vkládání obrázků
\usepackage{fancyvrb}       % vylepšené prostředí pro strojové písmo
\usepackage{indentfirst}    % zavede odsazení 1. odstavce kapitoly
\usepackage[numbers]{natbib}         % zajištuje možnost odkazovat na literaturu
			    % stylem AUTOR (ROK),~resp. AUTOR [ČÍSLO]
\usepackage[nottoc]{tocbibind} % zajistí přidání seznamu literatury,
                            % obrázků a~tabulek do obsahu
% \usepackage{icomma}         % inteligetní čárka v~matematickém módu
\usepackage{dcolumn}        % lepší zarovnání sloupců v~tabulkách
\usepackage{booktabs}       % lepší vodorovné linky v~tabulkách
\usepackage{paralist}       % lepší enumerate a~itemize
\usepackage{xcolor}         % barevná sazba
\usepackage{float}
\usepackage{ifthen}
\usepackage{cancel}
\usepackage{enumitem}
\usepackage{ragged2e}
\usepackage{needspace}		
\usepackage{wrapfig}		% plovoucí obrázky nalevo/napravo
\usepackage{pgfplots}
\usepackage{subcaption}
\usepackage{multicol}
\pgfplotsset{compat=1.15}
\usepackage{mathrsfs}
\usetikzlibrary{arrows}

\justifying

\graphicspath{{./images/}}

%%% Údaje o~práci

% Název práce v~jazyce práce (přesně podle zadání)
\def\NazevPrace{Fraktální geometrie pro (zdatné) amatéry}

% Název práce v~angličtině
\def\NazevPraceEN{\todo{Fractal geometry for (experienced) amateurs}}

% Jméno autora
\def\AutorPrace{Bc. David Weber}

% Rok odevzdání
\def\RokOdevzdani{\todo{2024}}

% Název katedry nebo ústavu, kde byla práce oficiálně zadána
% (dle Organizační struktury MFF UK, případně plný název pracoviště mimo MFF)
\def\Katedra{Katedra didaktiky matematiky}
\def\KatedraEN{Department of Mathematics Education}

% Jedná se o~katedru (department) nebo o~ústav (institute)?
\def\TypPracoviste{Katedra}
\def\TypPracovisteEN{Department}

% Vedoucí práce: Jméno a~příjmení s~tituly
\def\Vedouci{RNDr.~Martin~Rmoutil,~Ph.D.}

% Pracoviště vedoucího (opět dle Organizační struktury MFF)
\def\KatedraVedouciho{Katedra didaktiky matematiky}
\def\KatedraVedoucihoEN{Department of Mathematics Education}

% Studijní program a~obor
\def\StudijniProgram{Učitelství matematiky pro střední školy}
\def\StudijniObor{Učitelství matematiky pro střední školy se sdruženým studiem Učitelství informatiky pro střední školy}

% Nepovinné poděkování (vedoucímu práce, konzultantovi, tomu, kdo
% zapůjčil software, literaturu apod.)
\def\Podekovani{%
\todo{Doplnit poděkování}
\input{acknowledgements.txt}
}

% Abstrakt (doporučený rozsah cca 80-200 slov; nejedná se o~zadání práce)
\def\Abstrakt{%
\todo{Doplnit abstrakt (česky)}
\input{abstract_cz.txt}
}
\def\AbstraktEN{%
\todo{Doplnit abstrakt (anglicky)}
\input{abstract_en.txt}
}

% 3 až 5 klíčových slov (doporučeno), každé uzavřeno ve složených závorkách
\def\KlicovaSlova{%
{}, {klicova slova}
}
\def\KlicovaSlovaEN{%
{keywords}
}

%% Balíček hyperref, kterým jdou vyrábět klikací odkazy v~PDF,
%% ale hlavně ho používáme k~uložení metadat do PDF (včetně obsahu).
%% Většinu nastavítek přednastaví balíček pdfx.
\hypersetup{unicode}
\hypersetup{breaklinks=true}

%% Definice různých užitečných maker (viz popis uvnitř souboru)
%%% Tento soubor obsahuje definice různých užitečných maker a~prostředí %%%
%%% Další makra připisujte sem, ať nepřekáží v~ostatních souborech.     %%%

%%% Drobné úpravy stylu

% Tato makra přesvědčují mírně ošklivým trikem LaTeX, aby hlavičky kapitol
% sázel příčetněji a~nevynechával nad nimi spoustu místa. Směle ignorujte.
% \makeatletter
% \def\@makechapterhead#1{
%  {\parindent \z@ \raggedright \normalfont
%   \Huge\bfseries \thechapter. #1
%   \par\nobreak
%   \vskip 20\p@
% }}
% \def\@makeschapterhead#1{
%  {\parindent \z@ \raggedright \normalfont
%   \Huge\bfseries #1
%   \par\nobreak
%   \vskip 20\p@
% }}
% \makeatother

\setlength{\parskip}{0.3em}

% Toto makro definuje kapitolu, která není očíslovaná, ale je uvedena v~obsahu.
\def\chapwithtoc#1{
\chapter*{#1}
\addcontentsline{toc}{chapter}{#1}
}

% Trochu volnější nastavení dělení slov, než je default.
\lefthyphenmin=2
\righthyphenmin=2

% Zapne černé "slimáky" na koncích řádků, které přetekly, abychom si
% jich lépe všimli.
\overfullrule=1mm

%%% Makra pro definice, věty, tvrzení, příklady, ... (vyžaduje baliček amsthm)

\theoremstyle{plain}
\newtheorem{theorem}{Věta}[section]
\newtheorem{lemma}[theorem]{Lemma}
\newtheorem{proposition}[theorem]{Tvrzení}
\newtheorem{corollary}[theorem]{Důsledek}
\newtheorem*{proposition*}{Tvrzení}

\theoremstyle{definition}
\newtheorem{definition}[theorem]{Definice}
\newtheorem{example}[theorem]{Příklad}
\newtheorem{remark}[theorem]{Poznámka}
\newtheorem{convention}[theorem]{Úmluva}
\newtheorem{denoting}[theorem]{Značení}

%%% Prostředí pro důkazy

\renewenvironment{proof}[1][]{
  \par\medskip\noindent
  \textit{\ifthenelse{\equal{#1}{}}
  {Důkaz}
  {#1}}.
}{
\hspace*{\fill}$\qedsymbol$\par\medskip
}
\newenvironment{solution}[1][]{
  \par\medskip\noindent
  \textit{\ifthenelse{\equal{#1}{}}
  {Řešení}
  {#1}}.
}{
\hspace*{\fill}$\qedsymbol$\par\medskip
}

%%% Prostředí pro sazbu kódu, případně vstupu/výstupu počítačových
%%% programů. (Vyžaduje balíček fancyvrb -- fancy verbatim.)

\DefineVerbatimEnvironment{code}{Verbatim}{fontsize=\small, frame=single}

%%% Prostor reálných čísel, přirozených čísel, ...
\newcommand{\R}{\mathbb{R}}
\newcommand{\C}{\mathbb{C}}
\newcommand{\N}{\mathbb{N}}
\newcommand{\Q}{\mathbb{Q}}
\newcommand{\Z}{\mathbb{Z}}

%%% Užitečné operátory pro statistiku a~pravděpodobnost
\DeclareMathOperator{\pr}{\textsf{P}}
\DeclareMathOperator{\E}{\textsf{E}\,}
\DeclareMathOperator{\var}{\textrm{var}}
\DeclareMathOperator{\sd}{\textrm{sd}}

%%% Goniometrické a cyklometrické funkce
\DeclareMathOperator{\cotg}{cotg}
\DeclareMathOperator{\arccotg}{arccotg}
\DeclareMathOperator{\arctg}{arctg}
\DeclareMathOperator{\tg}{tg}

%%% Příkaz pro transpozici vektoru/matice
\newcommand{\T}[1]{#1^\top}

%%% Vychytávky pro matematiku
\newcommand{\goto}{\rightarrow}
\newcommand{\gotop}{\stackrel{P}{\longrightarrow}}
\newcommand{\maon}[1]{o(n^{#1})}
\newcommand{\abs}[1]{\left|{#1}\right|}
\newcommand{\dint}{\int_0^\tau\!\!\int_0^\tau}
\newcommand{\isqr}[1]{\frac{1}{\sqrt{#1}}}

\newcommand{\napprox}{\not\approx}
\newcommand{\norm}[1]{|| #1 ||}
\newcommand{\dotprod}[2]{\left(#1|#2\right)}
\newcommand{\compl}[1]{#1^\complement}
\newcommand{\map}[3]{#1:#2 \rightarrow #3}
\newcommand{\mapto}[3]{#1:#2 \mapsto #3}
\newcommand{\powset}[1]{\mathcal{P}\left(#1\right)}
\newcommand{\solutions}[1]{\left[K=\left\{#1\right\}\right]}
\newcommand{\set}[1]{\left\{#1\right\}}
\newcommand{\admid}{\;\middle\vert\;}
\newcommand{\sizeof}[1]{\left|#1\right|}
\newcommand{\dx}[1][]{
   \ifthenelse{\equal{#1}{}}%
      {\ensuremath{\,\mathrm{d}x}}%
      {\ensuremath{\,\mathrm{d}#1}}%
} % diferenciál

\newcommand{\ZF}{\textsf{ZF}} % Zermelova-Fraenkelova teorie množin
\newcommand{\PA}{\textsf{PA}} % Peanova aritmetika

% Čítače
\newcommand{\createcnt}[2][0]{\newcounter{#2}\setcounter{#2}{#1}}
\newcommand{\printcnt}[1]{\arabic{#1}}
\newcommand{\printnstepcnt}[1]{\stepcounter{#1}\arabic{#1}}

%%% Vychytávky pro tabulky
\newcommand{\pulrad}[1]{\raisebox{1.5ex}[0pt]{#1}}
\newcommand{\mc}[1]{\multicolumn{1}{c}{#1}}

% Celá jména
\newcommand{\name}[1]{\mbox{\textsc{#1}}}

% Logické operátory
\renewcommand{\implies}{\Longrightarrow}
\renewcommand{\impliedby}{\Longleftarrow}
\renewcommand{\iff}{\Longleftrightarrow}

% Cesty
\newcommand{\chapterpath}[1]{components/ch#1}
\newcommand{\sectionpath}[1]{components/ch#1/sections}
\newcommand{\literaturepath}{components/literature}
\newcommand{\appendixpath}{components/appendix}

% TODO
\newcommand{\todo}[1]{\textcolor{red}{(\noindent TODO: #1.)}}

% Rightsided note
\newcommand{\rightnote}[1]{\hspace*{\fill} $\triangleleft$ \textit{#1}}

% Images
\newcommand{\subfigwidth}{6cm}        % šířka okénka pro "podobrázky"
\newcommand{\wrappedfigwidth}{6cm}
\newcommand{\fullhd}{0.2}             % měřítko původního Full HD obrázku
\newcommand{\normalipe}{0.8}          % standardní měřítko IPE obrázku
\newcommand{\fractalscale}{0.3}       % měřítko obrázku softwarem vygenerovaného fraktálu

%% Titulní strana a~různé povinné informační strany
\begin{document}
%%% Titulní strana práce a~další povinné informační strany

%%% Titulní strana práce

\pagestyle{empty}
\hypersetup{pageanchor=false}

\begin{center}

\centerline{\mbox{\includegraphics[width=166mm]{components/images/logo-cs.pdf}}}

\vspace{-8mm}
\vfill

{\bf\Large DIPLOMOVÁ PRÁCE}

\vfill

{\LARGE\AutorPrace}

\vspace{15mm}

{\LARGE\bfseries\NazevPrace}

\vfill

\Katedra

\vfill

{
\centerline{\vbox{\halign{\hbox to 0.45\hsize{\hfil #}&\hskip 0.5em\parbox[t]{0.45\hsize}{\raggedright #}\cr
Vedoucí diplomové práce:&\Vedouci \cr
\noalign{\vspace{2mm}}
Studijní program:&\StudijniProgram \cr
\noalign{\vspace{2mm}}
Studijní obor:&\StudijniObor \cr
}}}}

\vfill

% Zde doplňte rok
Praha \RokOdevzdani

\end{center}

\newpage

%%% Následuje vevázaný list -- kopie podepsaného "Zadání bakalářské práce".
%%% Toto zadání NENÍ součástí elektronické verze práce, nescanovat.

%%% Strana s čestným prohlášením k~bakalářské práci

\openright
\hypersetup{pageanchor=true}
\pagestyle{plain}
\pagenumbering{roman}
\vglue 0pt plus 1fill

\noindent
Prohlašuji, že jsem tuto bakalářskou práci vypracoval(a) samostatně a~výhradně
s~použitím citovaných pramenů, literatury a~dalších odborných zdrojů.
Tato práce nebyla využita k~získání jiného nebo stejného titulu.

\medskip\noindent
Beru na~vědomí, že se na moji práci vztahují práva a~povinnosti vyplývající
ze zákona č. 121/2000 Sb., autorského zákona v~platném znění, zejména skutečnost,
že Univerzita Karlova má právo na~uzavření licenční smlouvy o~užití této
práce jako školního díla podle §60 odst. 1 autorského zákona.

\vspace{10mm}

\hbox{\hbox to 0.5\hsize{%
V \hbox to 6em{\dotfill} dne \hbox to 6em{\dotfill}
\hss}\hbox to 0.5\hsize{\dotfill\quad}}
\smallskip
\hbox{\hbox to 0.5\hsize{}\hbox to 0.5\hsize{\hfil Podpis autora\hfil}}

\vspace{20mm}
\newpage

%%% Poděkování

\openright

\noindent
\Podekovani

\newpage

%%% Povinná informační strana bakalářské práce

\openright

\vbox to 0.5\vsize{
\setlength\parindent{0mm}
\setlength\parskip{5mm}

Název práce:
\NazevPrace

Autor:
\AutorPrace

\TypPracoviste:
\Katedra

Vedoucí bakalářské práce:
\Vedouci, \KatedraVedouciho

Abstrakt:
\Abstrakt

Klíčová slova:
\KlicovaSlova

\vss}

\newpage

\nobreak\vbox to 0.49\vsize{
\setlength\parindent{0mm}
\setlength\parskip{5mm}

Title:
\NazevPraceEN

Author:
\AutorPrace

\TypPracovisteEN:
\KatedraEN

Supervisor:
\Vedouci, \KatedraVedoucihoEN

Abstract:
\AbstraktEN

Keywords:
\KlicovaSlovaEN

\vss}

\newpage

\openright
\pagestyle{plain}
\pagenumbering{arabic}
\setcounter{page}{1}


%%% Strana s automaticky generovaným obsahem bakalářské práce

\tableofcontents

%%% Jednotlivé kapitoly práce jsou pro přehlednost uloženy v~samostatných souborech
\include{\chapterpath{00}/ch00_predmluva.tex}
\include{\chapterpath{01}/ch01_uvod_do_fraktalu.tex}
\include{\chapterpath{02}/ch02_teorie_miry.tex}
\include{\chapterpath{03}/ch03_klasifikace_fraktalu.tex}
\include{\chapterpath{04}/ch04_generovani_fraktalu.tex}
\include{\chapterpath{05}/ch05_fraktaly_v_praxi.tex}

%%% Seznam použité literatury
%%% Seznam použité literatury (bibliografie)
%%%
%%% Pro vytváření bibliografie používáme bibTeX. Ten zpracovává
%%% citace v~textu (např. makro \cite{...}) a~vyhledává k nim literaturu
%%% v~souboru literatura.bib.
%%%
%%% Příkaz \bibliographystyle určuje, jakým stylem budou citovány odkazy
%%% v~textu. V závorce je název zvoleného souboru .bst. Styly plainnat
%%% a~unsrt jsou standardní součástí latexových distribucí. Styl czplainnat
%%% je dodáván s touto šablonou a~bibTeX ho hledá v~aktuálním adresáři.

% \bibliographystyle{czplainnat}    %% Autor (rok) s českými spojkami
% \bibliographystyle{plainnat}    %% Autor (rok) s anglickými spojkami
\bibliographystyle{unsrtnat}       %% [číslo]
% \bibliographystyle{alpha}

\renewcommand{\bibname}{Seznam použité literatury}

%%% Vytvoření seznamu literatury. Pozor, pokud jste necitovali ani jednu
%%% položku, seznam se automaticky vynechá.

\bibliography{\literaturepath/literature.bib}

%%% Kdybyste chtěli bibliografii vytvářet ručně (bez bibTeXu), lze to udělat
%%% následovně. V takovém případě se řiďte normou ISO 690 a~zvyklostmi v~oboru.

% \begin{thebibliography}{99}
%
% \bibitem{lamport94}
%   {\sc Lamport,} Leslie.
%   \emph{\LaTeX: A Document Preparation System}.
%   2. vydání.
%   Massachusetts: Addison Wesley, 1994.
%   ISBN 0-201-52983-1.
%
% \end{thebibliography}


%%% Obrázky v~bakalářské práci
%%% (pokud jich je malé množství, obvykle není třeba seznam uvádět)
\listoffigures

%%% Tabulky v~bakalářské práci (opět nemusí být nutné uvádět)
%%% U matematických prací může být lepší přemístit seznam tabulek na začátek práce.
% \listoftables

%%% Přílohy k~bakalářské práci, existují-li. Každá příloha musí být alespoň jednou
%%% odkazována z vlastního textu práce. Přílohy se číslují.
%%%
%%% Do tištěné verze se spíše hodí přílohy, které lze číst a~prohlížet (dodatečné
%%% tabulky a~grafy, různé textové doplňky, ukázky výstupů z počítačových programů,
%%% apod.). Do elektronické verze se hodí přílohy, které budou spíše používány
%%% v~elektronické podobě než čteny (zdrojové kódy programů, datové soubory,
%%% interaktivní grafy apod.). Elektronické přílohy se nahrávají do SISu a~lze
%%% je také do práce vložit na CD/DVD. Povolené formáty souborů specifikuje
%%% opatření rektora č. 72/2017.

\appendix

\openright
\end{document}
